\section{Stabilizer counterexamples}
\label{sec:stabilizer}

The X-cube model provides an exact control in which the protected rank grows
rapidly but the quantum geometry is fully determined by the tangent algebra.
On the three-torus its dependent stabilizer relations leave $6L-3$ logical
qubits and ground-space dimension $2^{6L-3}$
\cite{vijay2016xcube}.  Neither a large $D$ nor a moving projector is therefore
sufficient to produce random-matrix curvature.  Figure~\ref{fig:mechanisms}(d)
shows two deformations that separate these possibilities without allocating
the exponentially large ground-space matrix.

The calculation is performed in the binary symplectic representation.  Each
Pauli stabilizer is a row $(\bm x|\bm z)$ over $\mathbb F_2$, commutation is a
symplectic inner product, and row reduction determines the number of
independent constraints.  This proves the logical-qubit count without
diagonalizing a $2^{3L^3}$-dimensional Hamiltonian.  It also exposes the
relation responsible for the transported curvature: the ordered product of
the chosen local generators is reduced against the stabilizer row space,
rather than inferred from a small-system spectrum.

\subsection{Positive coefficient reweighting}

First vary the positive coefficients multiplying the commuting stabilizers
while keeping the Pauli generators fixed.  The simultaneous $+1$ eigenspace
does not depend on those coefficients as long as none crosses zero.  Hence the
projector is parameter independent,
\begin{equation}
 \partial_aP=0,
 \qquad X_a=0,
 \qquad g_{ab}=F_{ab}=0.
 \label{eq:xcube-flat-control}
\end{equation}
The registered coefficient reweighting curvature is exactly
\XCubeCoefficientCurvature.  This is stronger than a small numerical
residual: it follows from the common stabilizer eigenspace.  The deformation
changes excitation energies and may change the gap scale, yet it cannot
rotate the protected subspace.  Energy dependence outside the fiber and
geometry of the fiber are independent data.

This control also identifies the relevant failure boundary.  A coefficient
passing through zero removes one constraint and closes the corresponding
gap; a logical operator added within the ground space makes $PHP$ nonscalar
and splits the exact multiplet.  Neither operation belongs to the open,
constant-rank family of Eq.~\eqref{eq:xcube-flat-control}.

Generic local Pauli fields need not split the torus ground space at first
order because a detectable error maps it out of the code space.  Their
finite-size splitting appears only at the perturbative order at which a
noncontractible logical operator is assembled.  That exponentially or
high-order small splitting is physically important for code stability, but
it is distinct from exact degeneracy at the commuting-projector point.  Our
controls stay on paths for which the algebra certifies exact equality of all
ground energies.

\subsection{Local isospectral unitary transport}

The second deformation conjugates the full commuting-projector Hamiltonian
by local unitaries.  Let
\begin{equation}
 H(\lambda_1,\lambda_2)=UHU^\dagger,
 \qquad
 U=e^{-i\lambda_1G_1}e^{-i\lambda_2G_2}.
 \label{eq:xcube-unitary-transport}
\end{equation}
The spectrum, degeneracy, and gap are then preserved exactly, while
$P(\boldsymbol\lambda)=UPU^\dagger$ moves.  The registered local Pauli
generators are individually detectable, anticommute, and have an ordered
product equal to a stabilizer up to phase.  Their commutator therefore acts
centrally on the code space.  The curvature is nonzero but scalar,
\begin{equation}
 F_{12}=\XCubeTransportCurvatureEigenvalue\,P,
 \qquad
 D^{-1}\Tr\!\left[
   \bigl(F_{12}-D^{-1}\Tr(F_{12})P\bigr)^2
 \right]=\XCubeTransportConnectedVariance.
 \label{eq:xcube-scalar-curvature}
\end{equation}
The dense low-rank toy audit fixes the sign in
Eq.~\eqref{eq:xcube-scalar-curvature}; the production statement follows from
binary symplectic identities for every opened size.  Thus unitary transport
is exactly isospectral and geometrically nontrivial, but its non-Abelian
curvature contains no fluctuating eigenvalue sector.  Parallel transport
around an infinitesimal loop gives a common phase rather than relative
mixing within the code space.

Because the two generators anticommute, their ordered infinitesimal
transports do not cancel.  Because their product is a stabilizer, however,
the resulting commutator has the same action on every code state.  This is a
concrete example of noncommuting motion in the ambient Hilbert space whose
projection to the fiber lies in the center.  It rules out the converse
shortcut that any noncommuting deformation must generate a fluctuating
non-Abelian spectrum.

The scalar curvature is compatible with a nonzero metric: the projector can
move in the ambient Hilbert space even though the commutator of two accessible
motions becomes central after restriction to the fiber.  What vanishes is the
traceless curvature component used for level-repulsion and connected-variance
tests.  This distinction prevents a nonzero Berry phase from being mistaken
for stochastic non-Abelian mixing.

The connected variance vanishes because the response algebra collapses to
the center on the protected fiber.  There is no ensemble of fluctuating
curvature matrices to which a Jacobi or Gaussian comparison should be
applied.  The complete-covariance gate is not applicable to this exact
control; labeling it ``passed'' or ``failed'' would confuse deterministic
centrality with a statistical null.

Together, the two X-cube paths exclude two shortcuts in interpreting the
other models.  The flat path shows that exact degeneracy can coexist with no
projector response.  The unitary transport path shows that nonzero response
and curvature can coexist with zero connected variance.  Thus large exact
degeneracy is not sufficient for geometric chaos; one must also establish a
noncentral accessible response algebra and then test its statistics with an
appropriate sampling measure.
